% =========================================================
%  references.tex
%  Rujukan modul Unity.
% =========================================================

\clearpage
\phantomsection
\addcontentsline{toc}{section}{Rujukan}
\section*{Rujukan}

\begin{thebibliography}{9}
\setlength{\itemsep}{0.35em}

\bibitem{bukuajar}
Cecep Suwanda. \emph{Pengantar Pemrograman Game: Buku Ajar Berbasis
\istilah{Outcome-Based Education} (OBE)}. Universitas Bale Bandung, 2026.
Bab XII--XVI menjadi sumber konsep modul ini.

\bibitem{rps}
Cecep Suwanda. \emph{Rencana Pembelajaran Semester \NamaMK\ (\KodeMK)}.
Program Studi \NamaProdi, Fakultas \NamaFakultas, Universitas Bale Bandung,
Semester \SemesterMK\ \TahunAkademik.

\bibitem{unity}
Unity Technologies. \emph{Unity Manual}. \url{https://docs.unity3d.com/Manual/},
diakses 2026.

\bibitem{leveldesign}
Unity Technologies. \emph{Introduction to Level Design in Game Development
and in Unity}. Rujukan teknik, bukan sumber teks.

\end{thebibliography}
